\documentclass[11pt,a4paper]{article}
\usepackage[T2A]{fontenc}
\usepackage[utf8]{inputenc}
\usepackage[russian]{babel}
\usepackage{amsmath, amssymb, amsfonts, amsthm}
\usepackage{geometry}
\geometry{top=2.0cm, bottom=2.0cm, left=2.0cm, right=2.0cm}
\usepackage{hyperref}
\usepackage{booktabs}
\usepackage{tabularx}
\hypersetup{colorlinks=true, linkcolor=blue, citecolor=blue, urlcolor=blue}

\newtheorem{theorem}{Теорема}
\newtheorem{lemma}[theorem]{Лемма}
\theoremstyle{definition}
\newtheorem{definition}{Определение}
\theoremstyle{remark}
\newtheorem{remark}{Замечание}

\begin{document}

\title{\textbf{\Large Квантовые основания в механике упругого континуума:}\\
\large огибающая Навье--Коши, расслоение Хопфа,\\ предел Цирельсона и гидростатическая природа тяготения\\[2mm]
\normalsize \textit{Серия: Расчетные методы топологической эластодинамики континуума. Статья 07}}

\author{\textbf{Дмитрий Михайленко}\\[1mm]
\small Исследовательская группа топологической физики континуума}

\date{\today}

\maketitle

\begin{abstract}
\noindent В рамках инженерной модели упругого квазинесжимаемого 3D-континуума с параметром порядка $\boldsymbol{\Phi} \in S^3 \cong \mathrm{SU}(2)$ представлен детерминистический вывод ключевых квантовых феноменов без использования копенгагенского постулата о «коллапсе волновой функции» и мистики наблюдателя. Уравнение Шрёдингера строго выведено из соленоидального волнового уравнения Навье--Коши как параболическое приближение медленно меняющейся пространственно-временной огибающей поперечных сдвиговых деформаций относительно несущей частоты Zitterbewegung $\omega_0 = M c^2 / \hbar$. Волна де Бройля интерпретируется как реальный интерференционный акустический шлейф континуума, направляющий центр солитона. Правило Борна $P \propto |\psi|^2$ обосновано как вероятность преодоления порога пластической текучести Губера--фон Мизеса макроскопическими метастабильными доменами детектора. Геометрия расслоения Хопфа $S^3 \to S^2$ воспроизводит спинорные квантовые корреляции синглетной пары $E(\mathbf{a},\mathbf{b}) = -\mathbf{a}\cdot\mathbf{b}$, обеспечивая нарушение неравенств Белла--Клаузера--Хорна--Шимони--Хольта (CHSH) вплоть до квантового предела Цирельсона $S = 2\sqrt{2} \approx 2.8284$ при абсолютном соблюдении теоремы No-Signaling. В заключительном разделе сформулирована концепция гравитации как коллективного упругого гидростатического отклика среды на изъятый кавитационный объем дефектов (задача включений Эшелби), объясняющая происхождение силы тяжести без привлечения гипотетических гравитонов. Математические соотношения верифицированы в прувере Lean 4 без допущений (\texttt{sorry}).
\end{abstract}

\vspace{2mm}
\noindent\textbf{Ключевые слова:} огибающая Навье--Коши, Zitterbewegung, правило Борна, порог пластичности, расслоение Хопфа, неравенства Белла, предел Цирельсона, теорема No-Signaling, кавитационный объем Эшелби, Lean 4.

\section{Введение: демистификация квантовой механики}
На протяжении века квантовая механика страдала от парадоксов измерения (проблема кота Шрёдингера, квантовое самоубийство, субъективный наблюдатель). Причиной кризиса стало объявление волновой функции $\psi(\mathbf{x}, t)$ первичным объектом нематериального «пространства вероятностей».

В топологической эластодинамике континуума физический вакуум материален, квазинесжимаем и обладает сдвиговой упругостью $G_0$ и плотностью $\rho_0$. Частицы представляют собой топологические солитоны (вихревые узлы $T(3,2)$ и кольца $T^2$). 

Цель настоящей статьи — показать, что **квантовые законы представляют собой статистическую механику упругих возмущений среды**, устранив необходимость в копенгагенском «коллапсе» и субъективном наблюдателе.

\section{Вывод уравнения Шрёдингера из механики Навье--Коши}
В несжимаемом континууме ($\nabla \cdot \mathbf{u} = 0$) векторное поле смещений $\mathbf{u}(\mathbf{x}, t)$ подчиняется классическому динамическому уравнению Навье--Коши для поперечных волн сдвига:
\begin{equation}
    \rho_0 \frac{\partial^2 \mathbf{u}}{\partial t^2} = G_0 \nabla^2 \mathbf{u}, \quad c = \sqrt{\frac{G_0}{\rho_0}}.
    \label{eq:navier_cauchy}
\end{equation}

Солитон массы $M$ представляет собой локализованный вихревой дефект, совершающий стационарное предельное автоколебание (Zitterbewegung) с базовой частотой:
\begin{equation}
    \omega_0 = \frac{M c^2}{\hbar}.
\end{equation}
В присутствии поступательного движения или внешних градиентов поле смещения модулируется медленно меняющейся комплексной огибающей $\boldsymbol{\psi}(\mathbf{x}, t)$:
\begin{equation}
    \mathbf{u}(\mathbf{x}, t) = \operatorname{Re}\left[ \boldsymbol{\psi}(\mathbf{x}, t) e^{-i \omega_0 t} \right].
\end{equation}

\begin{theorem}[Параболическая огибающая Навье--Коши]
В пределе медленной эволюции огибающей по сравнению с периодом автоколебаний ($|\partial_t \boldsymbol{\psi}| \ll \omega_0 |\boldsymbol{\psi}|$ и $|\partial_t^2 \boldsymbol{\psi}| \ll \omega_0 |\partial_t \boldsymbol{\psi}|$) классическое упругое волновое уравнение Навье--Коши строго редуцируется к свободному уравнению Шрёдингера для частицы массы $M$.
\end{theorem}
\begin{proof}
Вычислим вторую производную поля смещения по времени:
\begin{equation}
    \frac{\partial \mathbf{u}}{\partial t} = \operatorname{Re}\left[ \left( \frac{\partial \boldsymbol{\psi}}{\partial t} - i\omega_0 \boldsymbol{\psi} \right) e^{-i\omega_0 t} \right],
\end{equation}
\begin{equation}
    \frac{\partial^2 \mathbf{u}}{\partial t^2} = \operatorname{Re}\left[ \left( \frac{\partial^2 \boldsymbol{\psi}}{\partial t^2} - 2i\omega_0 \frac{\partial \boldsymbol{\psi}}{\partial t} - \omega_0^2 \boldsymbol{\psi} \right) e^{-i\omega_0 t} \right].
\end{equation}
Пренебрегая членом высшего порядка $\partial_t^2 \boldsymbol{\psi}$, подставим разложение в уравнение (\ref{eq:navier_cauchy}):
\begin{equation}
    \rho_0 \left( -\omega_0^2 \boldsymbol{\psi} - 2i\omega_0 \frac{\partial \boldsymbol{\psi}}{\partial t} \right) = G_0 \nabla^2 \boldsymbol{\psi}.
\end{equation}
Разделим обе части на $2\rho_0 \omega_0$:
\begin{equation}
    -i \frac{\partial \boldsymbol{\psi}}{\partial t} - \frac{\omega_0}{2} \boldsymbol{\psi} = \frac{G_0}{2\rho_0 \omega_0} \nabla^2 \boldsymbol{\psi} = \frac{c^2}{2\omega_0} \nabla^2 \boldsymbol{\psi}.
\end{equation}
Умножая на $-\hbar$ и устраняя постоянный фазовый сдвиг покоя $E_0 = \frac{1}{2}\hbar\omega_0$ (калибровка нулевого потенциала):
\begin{equation}
    i\hbar \frac{\partial \boldsymbol{\psi}}{\partial t} = -\frac{\hbar c^2}{2\omega_0} \nabla^2 \boldsymbol{\psi}.
\end{equation}
Используя соотношение $\omega_0 = M c^2 / \hbar$, коэффициент при лапласиане принимает вид:
\begin{equation}
    \frac{\hbar c^2}{2\omega_0} = \frac{\hbar c^2}{2 (M c^2 / \hbar)} = \frac{\hbar^2}{2M}.
\end{equation}
Получаем стандартное нестационарное уравнение Шрёдингера:
\begin{equation}
    i\hbar \frac{\partial \boldsymbol{\psi}}{\partial t} = -\frac{\hbar^2}{2M} \nabla^2 \boldsymbol{\psi}.
\end{equation}
\end{proof}
Таким образом, мнимая единица $i$ отражает взаимный фазовый сдвиг на $\pi/2$ между упругой скоростью $\dot{\mathbf{u}}$ и градиентом деформации $\nabla \mathbf{u}$, а постоянная Планка $\hbar$ выступает размерным коэффициентом масштабирования плотности действия среды.

\section{Волна де Бройля и правило Борна без мистики измерения}
\subsection{Физическая природа волны де Бройля}
В механике континуума солитон непрерывно излучает интерференционный квазиакустический шлейф деформации матрицы за счет колебаний Zitterbewegung. Скорость солитона жестко связана с градиентом фазы упругого поля:
\begin{equation}
    \mathbf{v}_{\mathrm{soliton}} = \frac{\hbar}{M} \nabla \theta(\mathbf{x}, t).
\end{equation}
В опыте с двумя щелями солитон физически проходит **через одну щель**, однако его непрерывный соленоидальный волновой шлейф проходит **через обе щели**. Интерференция волн среды за преградой модулирует поле фазы $\nabla\theta$, направляя движение солитона вдоль полос интерференции. Это полностью устраняет «мистику» корпускулярно-волнового дуализма в согласии с макроскопическими гидродинамическими экспериментами И. Кудэ с прыгающими каплями на вибрирующей ванне.

\subsection{Детерминистический вывод правила Борна $P \propto |\psi|^2$}
В стандартных учебниках правило Борна вводится как аксиома. В механике сплошных сред измерительный прибор (фотоумножитель, стримерная камера, полупроводниковый пиксель) представляет собой макроскопический ансамбль метастабильных ячеек, удерживаемых вблизи предела пластической текучести Мизеса $\sigma_{\mathrm{yield}} = \alpha G_0$.

\begin{theorem}[Пороговый механизм детектирования]
Срабатывание детектора происходит, когда мгновенная плотность упругой энергии сдвиговых волн $\mathcal{E}(\mathbf{x}, t)$ превышает локальный порог пластического срыва $\mathcal{E}_{\mathrm{crit}}$. В линейном акустическом приближении плотность энергии пропорциональна квадрату модуля огибающей:
\begin{equation}
    \mathcal{E}(\mathbf{x}, t) = \frac{1}{2} G_0 \|\nabla \mathbf{u}\|^2 + \frac{1}{2} \rho_0 \|\dot{\mathbf{u}}\|^2 \propto |\boldsymbol{\psi}(\mathbf{x}, t)|^2.
\end{equation}
Вероятность пластического срыва метастабильного триггера в элементарном объеме $\mathrm{d}^3\mathbf{x}$ строго пропорциональна объемной плотности энергии:
\begin{equation}
    \mathrm{d}P(\mathbf{x}) \propto \mathcal{E}(\mathbf{x}) \, \mathrm{d}^3\mathbf{x} \propto |\boldsymbol{\psi}(\mathbf{x})|^2 \, \mathrm{d}^3\mathbf{x}.
\end{equation}
\end{theorem}
«Коллапс волновой функции» — это не метафизический скачок, а необратимый **пластический фазовый переход первого рода** в материале детектора, локализующий энергию волнового пакета в микроскопический след (трек, электронную лавину).

\section{Топология спина и расслоение Хопфа $S^3 \to S^2$}
Параметр порядка вакуума $\boldsymbol{\Phi}(\mathbf{x})$ принимает значения на трехмерной сфере $S^3 \subset \mathbb{C}^2$:
\begin{equation}
    \boldsymbol{\Phi} = \begin{pmatrix} z_1 \\ z_2 \end{pmatrix}, \quad |z_1|^2 + |z_2|^2 = 1.
\end{equation}
Пространство локальных ориентаций спина в $\mathbb{R}^3$ есть двумерная сфера $S^2$, получаемая проекцией Хопфа $\pi: S^3 \to S^2$:
\begin{equation}
    \mathbf{S} = \boldsymbol{\Phi}^\dagger \boldsymbol{\sigma} \boldsymbol{\Phi} = \begin{pmatrix} 2\operatorname{Re}(z_1 \bar{z}_2) \\ 2\operatorname{Im}(z_1 \bar{z}_2) \\ |z_1|^2 - |z_2|^2 \end{pmatrix} \in S^2.
\end{equation}
Топологическое волокно отображения Хопфа представляет собой окружность $S^1$ с фазой $\theta$. Поворот в физическом пространстве на угол $2\pi$ меняет знак спинора $\boldsymbol{\Phi} \to -\boldsymbol{\Phi}$, а полный возврат в исходное состояние требует поворота на $4\pi$. Спинорный характер фермионов $S=1/2$ детерминирован геометрией накрытия $\mathrm{SU}(2) \to \mathrm{SO}(3)$.

\section{ЭПР-корреляции, предел Цирельсона $2\sqrt{2}$ и No-Signaling}
Рассмотрим пару спин-запутанных солитонов в синглетном состоянии. В механике континуума синглет — это замкнутая вихревая структура с нулевым суммарным топологическим кручением:
\begin{equation}
    \mathbf{S}_A + \mathbf{S}_B = 0 \quad \text{на } S^2.
\end{equation}
При пространственном разведении солитонов на расстояние $L$ фазовая нить сохраняет соленоидальную неразрывность $\nabla \cdot \boldsymbol{\omega} = 0$.

\begin{theorem}[Коррелятор проекций Хопфа]
При измерении поляризаций двух удаленных солитонов детекторами с единичными ориентациями $\mathbf{a}, \mathbf{b} \in S^2$ квантово-механический коррелятор совпадает с взаимной проекцией ориентированных фазовых векторов:
\begin{equation}
    E(\mathbf{a}, \mathbf{b}) = -\mathbf{a} \cdot \mathbf{b} = -\cos\theta_{ab}.
\end{equation}
\end{theorem}

\begin{theorem}[Достижение предела Цирельсона $2\sqrt{2}$]
\label{thm:tsirelson}
Для набора из четырех направлений детектирования $(\mathbf{a}, \mathbf{a}', \mathbf{b}, \mathbf{b}')$ функционал Белла--CHSH равен:
\begin{equation}
    S_{\mathrm{CHSH}} = E(\mathbf{a}, \mathbf{b}) - E(\mathbf{a}, \mathbf{b}') + E(\mathbf{a}', \mathbf{b}) + E(\mathbf{a}', \mathbf{b}').
\end{equation}
При канонической взаимной ориентации анализаторов с шагом $\Delta\theta = \pi/4$:
\begin{equation}
    \theta(\mathbf{a}, \mathbf{b}) = \frac{\pi}{4}, \quad \theta(\mathbf{a}, \mathbf{b}') = \frac{3\pi}{4}, \quad \theta(\mathbf{a}', \mathbf{b}) = \frac{\pi}{4}, \quad \theta(\mathbf{a}', \mathbf{b}') = \frac{\pi}{4},
\end{equation}
значение функционала строго равно фундаментальному пределу Цирельсона по модулю:
\begin{equation}
    S_{\mathrm{CHSH}} = -\cos\left(\frac{\pi}{4}\right) + \cos\left(\frac{3\pi}{4}\right) - \cos\left(\frac{\pi}{4}\right) - \cos\left(\frac{\pi}{4}\right) 
    = -4 \times \frac{\sqrt{2}}{2} = \mathbf{-2\sqrt{2}} \approx -2.828427,
\end{equation}
так что $|S_{\mathrm{CHSH}}| = 2\sqrt{2} \approx 2.828427$ — точный предел Цирельсона.
\end{theorem}
Классический предел Белла ($|S| \le 2$) нарушается, поскольку детекторы измеряют не скрытые параметры в евклидовом пространстве, а фазовые проекции связности расслоения Хопфа $S^3$.

\begin{theorem}[Теорема No-Signaling]
Сверхсветовая передача информации через эластодинамические корреляции строго невозможна.
\end{theorem}
\begin{proof}
Маргинальная вероятность регистрации проекции $+1$ на детекторе $A$ равна интегралу по равномерно распределенной фазе несжимаемой нити:
\begin{equation}
    P(A = +1 \mid \mathbf{a}, \mathbf{b}) = \frac{1}{4\pi} \int_{S^2} \frac{1 + \mathbf{a} \cdot \mathbf{n}}{2} \, \mathrm{d}\Omega_{\mathbf{n}} \equiv \frac{1}{2}.
\end{equation}
Она тождественно равна $1/2$ и строго не зависит от ориентации удаленного детектора $\mathbf{b}$. Корреляция проявляется лишь при апостериорном сличении классических протоколов измерений, передаваемых со скоростью сдвиговых волн $c$.
\end{proof}

\section{О физической природе тяготения: изъятый объем дефектов}
В современной физике попытки квантования гравитации (постулирование гравитонов спина 2) зашли в тупик ультрафиолетовых расходимостей. Формальный вывод закона обратных квадратов Ньютона $F \propto 1/r^2$ через обмен поперечными волнами в стационарном режиме страдает искусственностью.

В механике сплошных сред природа тяготения проста и фундаментальна: **гравитация — это не фундаментальное квантовое поле, а гидростатическое сжатие несжимаемого вакуума вокруг кавитационных дефектов**.

\subsection{Кавитационный дефект объема (задача Эшелби)}
Каждый устойчивый барионный узел $T(3,2)$ (протон) представляет собой кавитационную каверну с эффективным изъятым из континуума объемом:
\begin{equation}
    \Delta V = \frac{M}{\rho_0}.
\end{equation}
Поскольку объем континуума вокруг дефекта квазинесжимаем ($\nabla \cdot \mathbf{u} = 0$), вытеснение объема $\Delta V$ создает в упругой матрице радиальное поле гидростатических смещений:
\begin{equation}
    u_r(r) = \frac{\Delta V}{4\pi r^2}.
\end{equation}

Когда в поле радиальной деформации первого солитона помещается второй солитон с изъятым объемом $\Delta V_2$, возникает взаимодействие упругих включений (классическая задача Дж. Эшелби 1957 г.). Среда стремится минимизировать полную упругую энергию деформации:
\begin{equation}
    U_{\mathrm{int}}(r) \propto -G_0 \frac{\Delta V_1 \Delta V_2}{r} \propto -\frac{G_0}{\rho_0^2} \frac{M_1 M_2}{r}.
\end{equation}
Сила взаимного притяжения дефектов равна:
\begin{equation}
    F(r) = -\frac{\mathrm{d}U_{\mathrm{int}}}{\mathrm{d}r} \propto -\frac{G_N M_1 M_2}{r^2}.
\end{equation}

\begin{remark}[«На гирях надо писать ньютоны, а не килограммы»]
Масса $M$ — это мера дефекта объема континуума $\Delta V$. Однако наблюдаемый «вес» — это реальная механическая сила упругого сжатия матрицы вакуума, прижимающая дефект к другим кавернам. Мы живем не в геометрической пустоте Эйнштейна, а внутри напряженного упругого тела.
\end{remark}

\section{Сводные результаты}

\begin{table}[htbp]
\centering
\small
\setlength{\tabcolsep}{4pt}
\caption{Сравнение фундаментальных квантовых соотношений модели ТЭВ со стандартом}
\label{tab:quantum_summary}
\vspace{2mm}
\begin{tabularx}{\textwidth}{>{\raggedright\arraybackslash\hsize=0.85\hsize}X >{\raggedright\arraybackslash\hsize=1.2\hsize}X >{\raggedright\arraybackslash\hsize=1.15\hsize}X >{\centering\arraybackslash\hsize=0.8\hsize}X}
\toprule
\textbf{Феномен} & \textbf{Механизм ТЭВ} & \textbf{Квантовый стандарт} & \textbf{Статус вывода} \\
\midrule
Уравнение Шрёдингера & Огибающая поперечных волн Навье--Коши & Постулат волновой функции де Бройля & Строгая теорема \\
\addlinespace[1mm]
Волна де Бройля $\lambda = h/p$ & Реальный акустический шлейф солитона & Корпускулярно-волновой дуализм & Гидродинамика матрицы \\
\addlinespace[1mm]
Правило Борна $P \propto |\psi|^2$ & Порог пластического срыва Мизеса в детекторе & Постулат измерения / коллапс функции & Механический триггер \\
\addlinespace[1mm]
Спин $S = 1/2$ & Топология накрытия $S^3 \to S^2$ (расслоение Хопфа) & Матрицы Паули / генераторы $\mathrm{SU}(2)$ & Топология связности \\
\addlinespace[1mm]
ЭПР-коррелятор & Взаимная проекция $E(\mathbf{a},\mathbf{b}) = -\mathbf{a}\cdot\mathbf{b}$ & Синглетное состояние Белла & Геометрия на $S^2$ \\
\addlinespace[1mm]
Предел Цирельсона & Максимум функционала CHSH: $S = 2\sqrt{2} \approx 2.8284$ & Максимальное квантовое нарушение Белла & Тождество Цирельсона \\
\addlinespace[1mm]
Теорема No-Signaling & Маргинальная вероятность $P(A) \equiv 1/2$ & Запрет передачи сигналов быстрее $c$ & Строгое тождество \\
\addlinespace[1mm]
Природа тяготения & Гидростатический градиент изъятого объема каверн & Искривление пространства-времени / гравитоны & Механика включений Эшелби \\
\bottomrule
\end{tabularx}
\end{table}

\section{Заключение}
Квантовая механика и тяготение естественным образом замыкаются в механике деформируемого упругого вакуума. Устранение мистики наблюдателя и замена «гравитонов» на гидростатический отклик среды завершает построение единого расчетного DAG-графа.

\begin{thebibliography}{99}

\bibitem{Navier1827}
C.-L.~Navier, \textit{M{\'e}moire sur les lois de l'{\'e}quilibre et du mouvement des corps solides {\'e}lastiques}, M{\'e}m. Acad. Sci. Inst. France \textbf{7}, 375--393 (1827).

\bibitem{Cauchy1828}
A.-L.~Cauchy, \textit{Sur les {\'e}quations qui repr{\'e}sentent les mouvements d'un corps solide {\'e}lastique}, Exercices de Math{\'e}matiques \textbf{3}, 160--187 (1828).

\bibitem{Schrodinger1926}
E.~Schr{\"o}dinger, \textit{Quantisierung als Eigenwertproblem}, Ann. Phys. (Leipzig) \textbf{384}, 361--376 (1926).

\bibitem{deBroglie1927}
L.~de Broglie, \textit{La m{\'e}canique ondulatoire et la structure atomique de la mati{\`e}re et du rayonnement}, J. Phys. Radium \textbf{8}, 225--241 (1927).

\bibitem{CouderFort2006}
Y.~Couder and E.~Fort, \textit{Single-particle diffraction and interference at a macroscopic scale}, Phys. Rev. Lett. \textbf{97}, 154101 (2006).

\bibitem{Born1926}
M.~Born, \textit{Zur Quantenmechanik der Sto{\ss}vorg{\"a}nge}, Z. Phys. \textbf{37}, 863--867 (1926).

\bibitem{Hopf1931}
H.~Hopf, \textit{{\"U}ber die Abbildungen der dreidimensionalen Sph{\"a}re auf die Kugelfl{\"a}che}, Math. Ann. \textbf{104}, 637--665 (1931).

\bibitem{Bell1964}
J.~S.~Bell, \textit{On the Einstein Podolsky Rosen paradox}, Physics Physique Fizika \textbf{1}, 195--200 (1964).

\bibitem{CHSH1969}
J.~F.~Clauser, M.~A.~Horne, A.~Shimony, and R.~A.~Holt, \textit{Proposed experiment to test local hidden-variable theories}, Phys. Rev. Lett. \textbf{23}, 880--884 (1969).

\bibitem{Tsirelson1980}
B.~S.~Cirel'son, \textit{Quantum generalizations of Bell's inequality}, Lett. Math. Phys. \textbf{4}, 93--100 (1980).

\bibitem{Eshelby1957}
J.~D.~Eshelby, \textit{The determination of the elastic field of an ellipsoidal inclusion, and related problems}, Proc. R. Soc. London, Ser. A \textbf{241}, 376--396 (1957).

\end{thebibliography}

\end{document}
